\section{A Typed Universal-Composability Execution Layer}
\label{sec:uc}

The UC subsystem instantiates the same exact-semantics discipline at the level
of interacting machines.  Its operational core is a typed, single-activation
FIFO kernel.  Protocols, ideal functionalities, environments, adversaries,
simulators, and networks are role-specific wrappers around exact interactive
machines; their public Boolean games are obtained only after running the
kernel and erasing its costs.

\begin{designprinciple}
The global scheduler owns routing, queue order, corruption, and kernel
bookkeeping.  An ITM owns only its address-indexed local state and its exact
handlers.  Neither side can manufacture a message whose payload, endpoints,
and connection capability disagree.
\end{designprinciple}

\subsection{Sessions, ports, and address-owned machines}

A session identifier \decl{SID} consists of a root tag and a list-valued child
path.  The operation \decl{SID.child} appends one tag, while
\decl{SID.Ancestor} is prefix inclusion restricted to an equal root.  A
machine address is the pair of a session identifier and a machine name.  This
representation makes independent top-level executions disjoint and retains
the complete nested-session path in every address.

The structure \decl{PortSchema} gives, for each address, direction, and payload
type, a type of ports.  A value of \decl{CanConnect} witnesses a permitted
connection between an output and an input port carrying the same payload
type.  It also determines a \decl{RoutingPolicy}: direct delivery or one of
the observable adversarial routes, with authenticated, forgeable, and
broadcast capabilities distinguished.  A separate \decl{CanSendAs}
capability authorizes a controller to claim a source address.  Accordingly,
\decl{Message}, \decl{Incoming}, and \decl{Emission} carry their payload type,
typed endpoints, and connection witness.  Ill-typed delivery is therefore
excluded when a message is constructed rather than detected later by the
interpreter.

An \decl{ITMFamily} is indexed by the security parameter and address.  It owns
dependent families \decl{State}, \decl{Leakage}, \decl{Erasure}, and
\decl{Output}, together with exact \decl{init}, \decl{activate},
\decl{applyErasure}, and \decl{leak} handlers.  An activation consumes either
a typed incoming message or \decl{resume}.  It returns a new local state and
exactly one \decl{LocalAction}: yield, emit, authorized send-as, erase, spawn,
request corruption, or output.  A computation that needs several emissions
must store a continuation and request later resume activations.  In
particular, broadcast permission never triggers hidden kernel fanout; a
broadcast component serializes its recipients as ordinary activations.

\subsection{A single FIFO and an explicit corruption transition}

The global \decl{Configuration} uses a dependent function as its local-state
store: the address determines the type of the optional stored value.  No
\decl{Any} or dynamic cast occurs.  The configuration also stores a list of
\decl{QueuedEvent}, the corrupted-address set, an optional terminal output,
and a typed audit trace.  Both ordinary activations and corruption requests
occupy the same FIFO; enqueue appends at the tail and \decl{Kernel.stepOne}
removes only the head.

The corruption policy is deliberately separate from erasure semantics.
\decl{CorruptionPolicy} specifies the admissible corrupted sets and the legal
one-address transition.  Its invariant inside a configuration states both
that the current set is admissible and that a corrupted address retains no
honest state.  The supplied \decl{incorruptible}, \decl{static}, and
\decl{dynamic} policies are restrictions of this same transition interface.

\Cref{fig:uc-fifo-corruption} shows the two event paths.  An honest activation
is lazily initialized, activated once, and allowed one action.  An activation
addressed to a corrupted target is redirected through the network-control
adapter.  A permitted corruption request reads the current state (initializing
a dormant machine if necessary), calls the exact leakage handler, removes the
honest state, inserts the target into the corrupted set, and enqueues a typed
leakage notification.  Subsequent send-as succeeds only when both the static
\decl{CanSendAs} witness exists and the claimed source belongs to this exact
configuration's corrupted set.  A rejected attempt is charged and audited but
cannot enqueue its forged emission.

\begin{figure}[t]
  \centering
  \resizebox{.92\linewidth}{!}{%
  \begin{tikzpicture}[node distance=6mm and 9mm]
    \node[rootbox] (queue) {global FIFO\\\(e_0 :: e_1 :: \cdots\)};
    \node[archbox, right=of queue] (dequeue) {charged\\dequeue};
    \node[archbox, right=of dequeue] (kind) {head event};

    \node[archbox, below=8mm of queue] (activation)
      {typed activation};
    \node[archbox, right=of activation] (honest)
      {state read/init\\exact \decl{activate}};
    \node[certbox, right=of honest] (action)
      {one \decl{LocalAction}\\append resulting event};

    \node[warnbox, below=7mm of activation] (request)
      {corruption request};
    \node[warnbox, right=of request] (policy)
      {check policy\\read/init state};
    \node[warnbox, right=of policy] (commit)
      {exact \decl{leak}\\remove state, mark corrupted};
    \node[certbox, right=of commit] (leakage)
      {typed leakage\\append adversarial activation};

    \draw[archarrow] (queue) -- (dequeue);
    \draw[archarrow] (dequeue) -- (kind);
    \draw[archarrow] (kind.south) |- (activation.east);
    \draw[archarrow] (activation) -- (honest);
    \draw[certarrow] (honest) -- (action);
    \draw[archarrow] (kind.south) |- (request.east);
    \draw[archarrow] (request) -- (policy);
    \draw[archarrow] (policy) -- (commit);
    \draw[certarrow] (commit) -- (leakage);
  \end{tikzpicture}}
  \caption{The authoritative scheduler consumes one FIFO event per kernel
    step.  Corruption is an independently queued transition, so leakage and
    state removal cannot be hidden inside the requesting activation.  Rejected
    requests leave the corrupted set unchanged.}
  \label{fig:uc-fifo-corruption}
\end{figure}

Erasure follows the same explicit chronology.  Processing an \decl{erase}
action calls \decl{applyErasure}, writes the resulting state, records the
request, and queues a resume.  A later corruption therefore leaks that stored
post-erasure state.  The kernel proves the ordering and state-removal
invariant; the meaning and adequacy of a particular erasure or leakage
function remain obligations of the installed ITM.

Kernel bookkeeping is itself costed.  \decl{KernelOperation} is a typed
signature covering dequeue, initialization, state access, routing, enqueue,
erasure, corruption, and termination, and \decl{KernelAlgebra} is its unique
exact handler.  The named \decl{KernelAlgebra.zero} handler is available when
an explicitly free bookkeeping model is intended; zero cost is not silently
built into the runner.

\subsection{Finite execution and the closed-world certificate chain}

The sole finite runner is
\[
  \decl{Kernel.runCosted} :
  \Nat \to \mathsf{Configuration} \to
  \mathsf{RandCosted}\;M\;\mathsf{ExecutionResult}.
\]
Its fuel counts scheduler steps only.  An \decl{ExecutionOutcome} is a typed
machine output, timeout, or deadlock; the final configuration and trace remain
in \decl{ExecutionResult}.  Fuel is not a checked cost budget.  At the public
Boolean boundary, \decl{ExecutionOutcome.toBool} maps timeout and deadlock to
$\false$, and \decl{Kernel.decisionDist} first erases the exact costs.  Thus a
failed certificate cannot change a branch or probability: certificates are
propositions about the runner, not inputs inspected by it.

The complexity proof is built in three stages.  First,
\decl{ComponentCostCertificate} bounds initialization, activation, erasure,
and leakage handlers, while independent \decl{OperationBounds} bound the
kernel algebra.  A \decl{StepCostCertificate} places all of these charges below
one security-parameter-dependent atom.  The current branch audit fixes a
conservative maximum of ten atomic charges; \decl{step_sound} follows the
actual \decl{stepOne} branches and pads shorter branches using nonnegativity.
Then \decl{runCosted_sound} proves the ordered repeated-step budget
\[
  c_{\mathrm{run}} \;\leq\;
  \mathsf{repeatActivationCost}_M(f,c_{\mathrm{step}}).
\]
Because the resource monoid need not be commutative, repetition and bind keep
the scheduler's left-to-right order.

Second, a \decl{FuelCertificate} supplies semantic evidence that the selected
activation limit has no timeout in exact support and that adding fuel leaves
the erased value distribution unchanged.  This is a genuine termination and
stability obligation; it is not inferred from a resource upper bound.  Third,
\decl{PPTExecutionCertificate} combines the step certificate, activation
limit, a natural-number bound on the measured step cost, polynomial proofs for
both functions, an independent operational admission of the exact closed
runner at the claimed runtime, and the fuel certificate.  Its exported runtime is
\[
  T(s)=f(s)\,T_{\mathrm{step}}(s),
\]
and \decl{NatMeasure.map_nsmul} connects the exact repeated cost to this
product.  This annotation proof does not generate operational admission; the
stored witness is reused by \decl{toPPTMachine}.  The resulting
\decl{PPTMachine} executes the original
\decl{Kernel.runCosted}; it does not execute a projected or reconstructed
runner.

For comparing two worlds, \decl{CertifiedWorldPair.commonFuel} is the
pointwise maximum of their activation limits.  Polynomial closure of maximum,
the two extended fuel certificates, and the theorems \decl{real_noTimeout},
\decl{ideal_noTimeout}, \decl{real_stable}, and \decl{ideal_stable} establish a
shared finite experiment.  In particular, changing the simulator cannot
change the real game merely by changing the selected common fuel, and changing
the adversary cannot change the ideal game for that reason; these independence
properties are proved from stability.

\subsection{Real and ideal worlds}

The global \decl{ClosedWorldAddress} is a disjoint sum of environment, system,
adversarial, and network addresses.  The function \decl{dispatchFamily}
selects the corresponding exact handler definitionally and preserves the
address-indexed local-state type.  \decl{composeReal} assembles
environment--protocol--adversary--network, whereas \decl{composeIdeal}
assembles the same role-separated interface with
environment--functionality--simulator--network; see
\Cref{fig:uc-worlds-contexts}.  \decl{Protocol} and
\decl{IdealFunctionality}, and likewise \decl{Adversary} and
\decl{Simulator}, are distinct structures even though each ultimately owns an
\decl{AddressedITM}.

An \decl{Environment} must expose an address-indexed
\decl{ULift} \decl{Bool} output.
The closed-world decision function accepts that output only when its source is
in the environment summand.  Outputs from the protocol, functionality,
adversarial role, or network are mapped to $\false$ rather than being mistaken
for the distinguisher's bit.  A real/ideal certified pair additionally shares
the environment and corruption policy, and an \decl{Experiment} fixes the
same port schema, network, protocol, and ideal functionality used by all
quantified roles.

Computational UC emulation is then stated with the standard dependency order:
\begin{equation}
  \forall\,\mathcal{A}\in\mathsf{PPT},\quad
  \exists\,\mathcal{S}\in\mathsf{PPT},\quad
  \forall\,\mathcal{Z}\in\mathsf{PPT},\quad
  \mathsf{Real}_{\Pi,\mathcal{A},\mathcal{Z}}
    \;\approx\;
  \mathsf{Ideal}_{\mathcal{F},\mathcal{S},\mathcal{Z}}.
  \label{eq:uc-emulation}
\end{equation}
Here each quantified role carries a \decl{PPTAddressedITMCertificate} for the
experiment's explicit cost model (M) and \decl{NatMeasure}; the complete
role certificate includes operational admission of the exact handlers and
claimed runtime.  Complete closed-world builders additionally provide an
execution certificate, with its own admission for the composed runner, that
accounts for the protocol/functionality and network.  The relation
\decl{Indistinguishable} compares the two cost-erased Boolean games.
\decl{PerfectUCEmulates} replaces indistinguishability by pointwise equality,
and \decl{PerfectUCEmulates.ucEmulates} derives the computational statement by
showing that the advantage is identically zero.

\begin{figure}[t]
  \centering
  \resizebox{\linewidth}{!}{%
  \begin{tikzpicture}[node distance=4mm and 5mm]
    \node[rootbox] (zr) {environment\\\(\mathcal Z\)};
    \node[archbox, right=of zr] (pi) {protocol\\\(\Pi\)};
    \node[archbox, right=of pi] (ar) {adversary\\\(\mathcal A\)};
    \node[archbox, right=of ar] (nr) {network\\\(\mathcal N\)};
    \node[certbox, right=9mm of nr] (kr) {exact FIFO kernel\\real run};

    \node[rootbox, below=9mm of zr] (zi) {same environment\\\(\mathcal Z\)};
    \node[archbox, right=of zi] (fun) {functionality\\\(\mathcal F\)};
    \node[archbox, right=of fun] (sim) {simulator\\\(\mathcal S\)};
    \node[archbox, right=of sim] (ni) {same network\\\(\mathcal N\)};
    \node[certbox, right=9mm of ni] (ki) {exact FIFO kernel\\ideal run};

    \node[warnbox, below=10mm of fun, xshift=10mm] (context)
      {typed context \(C[\,\cdot\,]\)\\shared system hole and role transforms};
    \node[certbox, right=8mm of context] (simulation)
      {one-step \decl{KernelSimulation}\\after exact-cost erasure};
    \node[rootbox, right=8mm of simulation] (compose)
      {operational equality\\and \decl{uc_compose}};

    \draw[archarrow] (zr) -- (pi);
    \draw[archarrow] (pi) -- (ar);
    \draw[archarrow] (ar) -- (nr);
    \draw[certarrow] (nr) -- (kr);
    \draw[archarrow] (zi) -- (fun);
    \draw[archarrow] (fun) -- (sim);
    \draw[archarrow] (sim) -- (ni);
    \draw[certarrow] (ni) -- (ki);
    \draw[erasearrow] (kr) -- node[right,font=\scriptsize]{shared fuel} (ki);
    \draw[archarrow] (context) -- (simulation);
    \draw[certarrow] (simulation) -- (compose);
    \draw[archarrow] (pi.south) -- (context.north west);
    \draw[archarrow] (fun.south) -- (context.north east);
  \end{tikzpicture}}
  \caption{Real and ideal worlds are concrete four-component executions.  A
    context transforms both system roles through one typed hole, while its
    operational premise is a commuting single-step kernel simulation.}
  \label{fig:uc-worlds-contexts}
\end{figure}

\subsection{Operational contexts and composition}

Universal composition is phrased over \decl{ExecutableExperiment}, whose
real and ideal worlds are structurally assembled from exact components and
their certificates.  A \decl{SystemHole} is one transformation of a
system-owned \decl{AddressedITM}; the same \decl{fill} produces both the outer
protocol and outer functionality.  \decl{ContextBuilder} also transforms the
shared policy and network and must construct the corresponding real and ideal
execution data.  Hence the outer experiment cannot be an unrelated
preassembled game attached to a nominal hole.

Renaming is similarly explicit.  \decl{WorldRenaming} consists of four
role-preserving injective address maps, and \decl{PortTransport} maps
activations, emissions, and send-as capabilities while preserving targets and
routing policy.  A \decl{KernelSimulation} then maps dependent states,
leakage, erasures, outputs, configurations, queues, corrupted sets, and traces.
Its central premise is a commuting \decl{Kernel.stepOne} square after exact
cost erasure, together with initial-configuration, decision, policy, and
network-adapter laws.  Induction yields finite-run and Boolean-decision
commutation.

A \decl{Context} additionally provides PPT-preserving transformations
\decl{contextAdversary} and \decl{contextEnvironment}, plus an
environment-independent \decl{plugSimulator}.  Its real and ideal premises are
simulations between the actual certified worlds produced by the builder and
the transformed inner roles.  Identity and sequential composition compose
these single-step simulations; plugging is definitionally unital and
associative for structurally executable experiments.  Theorems
\decl{real_operational} and \decl{ideal_operational} turn the simulations into
equalities of the compared Boolean games.  Finally, \decl{uc_compose} applies
the inner simulator to the context-transformed adversary and lifts it with
\decl{plugSimulator}, preserving the quantifier order in
\Cref{eq:uc-emulation}.

\begin{scopeboundary}
The current context interface is intentionally endomorphic: it keeps one
global port schema and the same four local address types, up to
role-preserving injective endomaps.  Its simulation theorem preserves erased
operational behavior, not exact path costs.  The context must construct new
PPT-certified executions; polynomial closure is enforced by these output
types rather than inferred from an arbitrary host-language function.
\end{scopeboundary}

\subsection{The layered-MPC bridge}

The \decl{Layered} namespace connects the generic kernel to a concrete family
shape without postulating a separate execution semantics.  Parameters record
the number of parties per layer, the number of layers, and a per-cell
corruption threshold.  A party identifier contains both layer and position,
and its system address additionally contains the full \decl{SID}.  The
predicate \decl{Corruption.Eligible} permits only party addresses and enforces
the threshold independently for every exact session--layer pair.  Thus a child
session or a different root cannot consume another cell's corruption budget.

A \decl{PartyStep} is an actual exact ITM, not a detached transition
relation.  \decl{SystemComponents} installs party machines, an explicit
broadcast manager, an explicit corruption manager, and typed boundary
machines; its total address dispatcher becomes the real \decl{Protocol}.
Likewise, \decl{MPCFunctionality} contains an actual addressed ITM and becomes
an \decl{IdealFunctionality}.  \decl{ExecutableLayered} supplies the shared
layered corruption policy, both initial configurations, and both
\decl{PPTExecutionCertificate}s, and
\decl{toExecutableExperiment} places these components directly into the
generic real/ideal and context infrastructure.

\begin{scopeboundary}
This bridge establishes executable typed wiring and certificate obligations;
it is not by itself a security proof for a particular layered MPC protocol.
Concrete party algorithms, broadcast behavior, ideal functionality behavior,
simulator construction, and the simulations required by \decl{uc_compose}
must still be supplied.  Likewise, the finite runner models activation-bounded
executions; liveness beyond a provided \decl{FuelCertificate} is outside the
present kernel theorem.
\end{scopeboundary}
